% Shared exercise chunk: l1-bounded-attention
\begin{exercisechunk}
\item \textbf{Bounded scores and attention concentration.}
\label[exercise]{ex:bounded-attention}
Let $T\ge2$ be an integer and let $a,b\in\mathbb R$ satisfy $a\le b$.
Write $R=b-a$ for the score range. For a score vector $s\in[a,b]^T$,
define the probability vector $\alpha(s)$ by
\[
\alpha_j(s)=\frac{e^{s_j}}{\sum_{r=1}^T e^{s_r}},\qquad j\in[T].
\]
\begin{enumerate}[(a)]
\item
Compute $\sup_{s\in[a,b]^T}\max_{j\in[T]}\alpha_j(s)$ and give a score
vector attaining this value.
\item \emph{Optional extension.}
Writing $H(\alpha)=-\sum_{j=1}^T\alpha_j\log\alpha_j$ for the
\emph{entropy}, show that for every $s\in[a,b]^T$,
\[
\log T-R\le H(\alpha(s))\le\log T.
\]
Deduce that, for fixed $a,b$, $H(\alpha(s))/\log T\to1$ uniformly over
$s\in[a,b]^T$ as $T\to\infty$. This concerns entropy relative to its maximum;
it does not imply $\sum_{j=1}^T|\alpha_j(s)-1/T|\to0$.
\item For a fixed $c\in(0,1)$, determine the smallest $R\ge0$ for which
some $s\in[a,a+R]^T$ satisfies $\max_{j\in[T]}\alpha_j(s)\ge c$.
Give such a score vector at this minimum range.
Describe how this minimum range grows with $T$ for fixed $c$, and deduce
whether a score range independent of $T$ can maintain $\max_j\alpha_j(s)\ge c$
as $T\to\infty$.
\end{enumerate}
\end{exercisechunk}
